\paragraph{Conjugate-coordinate class times a fuzzy factor: the kinetic matrix of the static scalar singlets, the combination selected by the source and the pole of the neutral singlet (conditional; bounded no-go).}
This paragraph supplements the preceding one, in which the kinetic matrix of the singlets \(T\) and \(\mathsf PT\), a matter coupling and the combination that the static channel selects were not computed. Keep its setting and notation, with \(y_0=(y_0^0,0,0,0)\), so that \(\mathsf P^a{}_b=\delta^a_0\delta^0_b\), static \(q=k\hat q\) with \(q\cdot y_0=0\), and the wave \(A^a=\epsilon_1(y)\,T^a+\epsilon_y(y)\,(\mathsf PT)^a\) with independent profiles, for which \(\sigma_1=-(4\epsilon_1+\epsilon_y)\); the tree density is \((\rho_M/2g^2)\langle A\cdot L(A)\rangle\), the normalisation of the preceding paragraph.

\emph{Kinetic matrix.} For a function \(f\) of the base functions, exactly, \(L(fA)=fL(A)+\partial_\mu f\,B^\mu(A)+\partial_\mu\partial_\nu f\,R^{\mu\nu}(A)\), with \(B^\mu(A)=L(y^\mu A)-y^\mu L(A)\), \(R^{\mu\nu}(A)_a=-m^2\bigl(\eta^{\mu\nu}A_a-\tfrac12(\delta^\mu_aA^\nu+\delta^\nu_aA^\mu)\bigr)\) and no term of third order. By the localisation identity the part of the tree form symmetric under the exchange of the two profiles is that of waves with a common profile: its non-derivative matrix is \(\tfrac12(\langle w_i\cdot L(w_j)\rangle+\langle w_j\cdot L(w_i)\rangle)\), where on the fibre over \(y_0\) \(\langle T\cdot L(\mathsf PT)\rangle=2(\ell^2+m^2)\) and \(\langle\mathsf PT\cdot L(T)\rangle=0\), and its derivative part is the leading symbol. The antisymmetric part is of first order, with coefficients \(\langle w_i\cdot B^\mu(w_j)\rangle\) and \(\partial_\mu\langle T\cdot R^{\mu\nu}(\mathsf PT)\rangle\), all along \(y_0\) at \(y_0\), so it vanishes for static \(q\). At the frozen base point the tree form of the static wave is therefore
\[
\frac{\rho_M}{2g^2}(\ell^2+m^2)\Bigl[\bigl(6+\tfrac23m^2k^2\bigr)\epsilon_1^2+2\epsilon_1\epsilon_y\Bigr],
\]
so the kinetic matrix is \(\tfrac23m^2k^2(\ell^2+m^2)\rho_M/2g^2\) on \(\epsilon_1\) and zero in the other entries, since \(\mathsf PT\) vanishes on the fibre over \(y_0\) (exact Lie--Poisson algebra for both signs \(s\), fibre averages by an exact quadrature).

\emph{Source.} For \(X=(1+\kappa)T\) with constant \(\kappa=\epsilon_1+\epsilon_y\mathsf P\) the frame gives \(\delta\gamma^{\mu\nu}=m^2(\kappa^{\mu\nu}+\kappa^{\nu\mu})\). The part of the frame change with the derivatives of the profiles contributes to \(h_\beta^{00}\) neither at a point of the fibre over \(y_0\), where \(T^0=0\), nor through the fibre average of \(t\), since \(\langle T_\mu M^{\mu\rho}\rangle\) is along \(y_0\). For dust coupled to \(h_\beta\) as in the paragraph on the response of the scale of \(K\) the source is therefore \(\varrho_0h_\beta^{00}\) with
\[
h_\beta^{00}=2m^2\bigl[(2\beta_K-1)\epsilon_1+(\tfrac12\beta_K-1)\epsilon_y\bigr].
\]
On the channel along \(y^\mu y^\nu\), \(\epsilon_1=0\), this is the charge \(1-\tfrac12\beta_K\) of that paragraph, and with \(\epsilon_1=0\) the forms reproduce its response at large \(\tau\), \(-L\varrho_0^2(1-\tfrac12\beta_K)^2/(q^2+\mu)\) with \(\mu\) defined below. On the \(\sigma_1\)-neutral singlet \((\mathsf P_\perp-3\mathsf P)T\) the charge is \(3\) for every \(\beta_K\), so for (S1) and for (S2). Since the source acts on that singlet and the tree form couples \(\epsilon_y\) to \(\epsilon_1\), the static equations do not keep \(\epsilon_1=0\): the channel along \(y^\mu y^\nu\) alone does not solve them, and the neutral singlet is eliminated by its own equation.

\emph{Selected combination and the constant at \(q=0\).} Take as hypothesis that the one-loop and induced terms depend on the two singlets through \(\sigma_1\) alone, with the values on the channel along \(y^\mu y^\nu\) of the identity of the paragraph on the tadpole of the static scalar channel. In its normalisation the static form is then \((q^2+\mu)\sigma_1^2\) plus the tree form, with \(\mu=(L/2m^4)\sum_i(C_i-T_i)+q_*^2|_\chi\) summed over the terms other than the tree terms. Write \(r=(L/2m^4)(\rho_M/2g^2)(\ell^2+m^2)=\vartheta\varpi/m^2\). At fixed \(\sigma_1\) the neutral direction is \(\epsilon_1\), and eliminating it exactly turns the denominator of the first term of the vacuum-term display into
\[
q^2+\mu+\frac{r}{2-\tfrac23m^2k^2},
\]
whose value at \(q=0\) adds \(\vartheta\varpi/2m^2\), the stationary value of the preceding paragraph, for every \(\beta_K\). At \(q=0\) the equations select \(\epsilon_1=-\tfrac12\sigma_1\), \(\epsilon_y=\sigma_1\), shifted along the neutral singlet by a multiple of the source that gives a term of the response without a pole. On this combination \(\delta\gamma^{00}=\delta\gamma^{11}=\delta\gamma^{22}=\delta\gamma^{33}=-m^2\sigma_1\) in the rest frame of \(y_0\), so, read in \(h_\beta\), the ratio of the two potentials on it is \((1-\beta_K)/(1+\beta_K)\), and after the elimination the source acts on \(\sigma_1\) with charge \(-\tfrac12(1+\beta_K)\), in the sign convention in which the channel along \(y^\mu y^\nu\) has \(1-\tfrac12\beta_K\). Against the exchange of \(\chi\), with the other entries of \(\mu\) set to zero, the pole of \(\sigma_1\) has the Yukawa sign exactly when \(\tfrac14(\rho_M/g^2)(\ell^2+m^2)>U_{\mathrm{loop}}'^2/2U_{\chi\chi}\). The induced coefficient \(A\) drops out of this condition, which involves the supplied \(\rho_M/g^2\) and, through \(U'_{\mathrm{loop}}\) and \(U_{\chi\chi}\), the tree data of \(K\).

\emph{The pole of the neutral singlet.} The value \(r/2\) is obtained at a maximum over \(\epsilon_1\), whose own coefficient at fixed \(\sigma_1\) is \(-2r\). At \(q=0\) the static form has determinant \(-r(2\mu+r)\). So, with the tree kinetic term above and the induced term \(q^2\sigma_1^2\), or with any positive-definite kinetic matrix of the two singlets, wherever the pole of \(\sigma_1\) has the Yukawa sign (\(\mu+r/2>0\)) the static response has exactly one further pole, at \(q^2>0\), of the tachyonic sign. At \(\mu=0\) the two poles are at \(m^2q^2=\tfrac32\bigl(1\pm\sqrt{1+\tfrac23\vartheta\varpi}\bigr)\), and \(|q|\ell\gg1\) holds at them only for \(\vartheta\varpi\) large against \(\sinh^4\tau\). More generally, take loop entries \(2b\,\sigma_1\epsilon_1+\nu\epsilon_1^2\) on the singlet sector with the value \(\mu\) on the channel along \(y^\mu y^\nu\) held. The constant after the elimination is then \(\mu-(b-r)^2/(\nu-2r)\), which exceeds \(\mu\) only where \(\nu<2r\). For a positive-definite kinetic matrix, a static response without a pole at \(q^2\ge0\) needs the form at \(q=0\) to be positive definite, hence \(\mu>0\) and a constant at most \(\mu\). Within these hypotheses the singlet sector therefore moves the tachyonic sign of the static scalar pole into the neutral singlet and does not remove it: a static scalar channel whose poles all have the Yukawa sign needs the same positive entries as the channel along \(y^\mu y^\nu\) alone. A form of Einstein--Hilbert type does not satisfy the hypothesis of a positive-definite kinetic matrix: on these metric data the Fierz--Pauli form is \(4m^4q^2(3\epsilon_1^2+2\epsilon_1\epsilon_y)\), which is indefinite and vanishes on the channel along \(y^\mu y^\nu\). The induced kinetic matrix of the singlets is not computed.

These statements are for the setting of the preceding paragraph, the frozen base point \(y_0\) with static \(q\), the recorded dust coupling with exponent \(\beta_K\), and, where stated, the hypothesis that the loop and induced terms depend on the singlets through \(\sigma_1\) alone. The operator identities, the kinetic matrix, the source and the forms are exact for these hypotheses, and the fibre averages are exact quadratures. They do not compute the induced or loop kinetic matrix of the singlets, the loop entries \(b\) and \(\nu\) or the entries of \(\mu\), do not derive the matter coupling, do not determine the sign of the static scalar pole, and do not identify \(\sigma_1\), \(t\) or the singlets with quantities of the emergent metric of Chapter~\ref{ch:vierbein-emergence}, so no Cassini or Newtonian comparison follows. Problems~C, F and~N remain open and PROGRAMME-CONDITIONAL.

